\begin{statement}{C.2.UN.6}{Proved}{Zero-Euler product fillings with essential Kodaira-dimension-one components}
Put $q=e^{-2\pi}$ and $E=\C^*/q^{\Z}\simeq\C/(\Z+\mathrm i\Z)$.
For every integer $m\ge2$, the fixed punctured product
$(E\times E)\times\Delta^*$ has a proper filling
$f_m:U_m\to\Delta$ with smooth total space and full central fibre
\[
                    (U_m)_0=P_m+mQ_m.             \tag{C.ER.1}
\]
Here $Q_m$ is a primary Hopf surface, and $P_m$ is a smooth compact
non-K\"ahler elliptic surface over $E$ with precisely one multiple
fibre $mJ$. The surfaces meet transversely along the smooth elliptic
curve $J$. One has
\[
 \begin{gathered}
 e((U_m)_0{}_{\rm red})=0,\qquad H_1(P_m,\Z)=\Z^3,\\
 K_{P_m}\simeq\OO_{P_m}((m-1)J),\qquad \kappa(P_m)=1,\\
 h^0(P_m,K_{P_m}^{\otimes r})=
 \begin{cases}
 1,&d_r=0,\\
 d_r,&d_r>0,
 \end{cases}
 \qquad d_r=\left\lfloor\frac{r(m-1)}m\right\rfloor,
 \quad r\ge1.
 \end{gathered}                                   \tag{C.ER.2}
\]
Neither $P_m$ nor $Q_m$ can be contracted by a bimeromorphic map
between smooth threefold fillings. The $U_m$ are pairwise
non-bimeromorphic over the disc. After every finite base change and
resolution of its normalization, a central component remains
birational to $P_m$. Thus a properly elliptic component of Kodaira
dimension one is forced to survive in these examples, including
after finite base change. All these fillings have compact smooth
threefold realizations over $\PP^1$.
\end{statement}
\begin{proof}
We give the construction as an actual hypersurface and a free
quotient. Let
\[
 V=\operatorname{Bl}_{(0,0)}(E\times\Delta),\quad
 R=\operatorname{Exc}(V\to E\times\Delta),\quad
 L=\OO_V(R),\quad \sigma\in H^0(V,L)
                                                        \tag{C.ER.3}
\]
with $\sigma$ the canonical section. Define
\[
 W_m=\{(v,\ell,a):v\in V,\ \ell\in L_v^\times,\ a\in\C,
                         \ \sigma(v)=a^m\ell\}.
                                                        \tag{C.ER.4}
\]
The generator of a $\Z$-action is
\[
             (v,\ell,a)\longmapsto(v,q^m\ell,q^{-1}a),
 \qquad U_m=W_m/\Z.                               \tag{C.ER.5}
\]
The action preserves (C.ER.4). Near $R$, choose a holomorphic frame
of $L$ in which $\sigma=xe$ and write $\ell=be$, $b\ne0$.
Then (C.ER.4) is $x=a^m b$ and eliminates the coordinate $x$;
thus $W_m$ is smooth. It is also smooth off $R$, where $a\ne0$
and $\ell=\sigma/a^m$. A Hermitian norm on $L$ is multiplied by
$q^m$ under (C.ER.5). It is bounded above and bounded away from zero
on each compact subset of $W_m$, proving proper discontinuity.
No nonzero power fixes a point, since $\ell\ne0$. Consequently
$U_m$ is a smooth complex threefold.

We verify properness and the local maps explicitly. With $y$ a
second base coordinate, the map
\[
 (z,w,y)\longmapsto
 (x=z^m,y;\ b=w^m,\ a=z/w),\qquad w\in\C^*,
                                                        \tag{C.ER.6}
\]
identifies the local $W_m$ with the quotient by
$(z,w,y)\mapsto(\zeta z,\zeta w,y)$, $\zeta^m=1$.
It is onto because an $m$th root of $b$ determines $w$ and then
$z=aw$; its fibres are exactly the displayed orbits. The action is
free. The generator (C.ER.5) lifts to $(z,w,y)\mapsto(z,qw,y)$.
Hence the local projection $\rho_m:U_m\to V$ is exactly
\[
 (\Delta_z\times E\times\Delta_y)/\boldsymbol\mu_m
        \longrightarrow\Delta_x\times\Delta_y,
 \quad [z,[w],y]\longmapsto(z^m,y),
 \quad \zeta:(z,[w],y)\longmapsto(\zeta z,[\zeta w],y).
                                                        \tag{C.ER.7}
\]
The last finite action is free: $[\zeta w]=[w]$ would imply
$\zeta=q^j$ for an integer $j$, forcing $j=0$ and $\zeta=1$.
The map before this finite quotient is the product of a finite
disc map with a compact elliptic curve. It is proper, so (C.ER.7)
is proper. These charts prove that $\rho_m$ is proper everywhere.
Its fibres are connected elliptic curves, including the reduced
fibre over $R$. Thus composition with $V\to\Delta$ is proper
with connected fibres.

The comparison with the punctured product is particularly simple.
For $v\notin R$ and $c\in\C^*$ set
\[
       (v,[c])\longmapsto[(v,\sigma(v)c^m,c^{-1})].
                                                        \tag{C.ER.8}
\]
Changing $c$ to $qc$ changes the representative by (C.ER.5).
Conversely $[c]=[a^{-1}]$ is its inverse. Thus (C.ER.8) is an
actual isomorphism $(V\setminus R)\times E\simeq
\rho_m^{-1}(V\setminus R)$, for every $m$, and over $\Delta^*$
it gives the specified fixed product $(E\times E)\times\Delta^*$.

Let $E_0\simeq E$ be the strict transform of $E\times\{0\}$.
It meets $R\simeq\PP^1$ at one point, and
$L|_{E_0}=\OO_E([0])$, $L|_R=\OO_{\PP^1}(-1)$.
The inverse image of $R$, with its reduction, is $a=0$ in
(C.ER.4), so it is
\[
 Q_m=L|_R^\times/(q^m)^{\Z}
     \simeq(\C^2\setminus\{0\})/\langle q^m\operatorname{id}\rangle.
                                                        \tag{C.ER.9}
\]
The first isomorphism sends the nonzero tautological vector to
the same vector in $\C^2\setminus\{0\}$. It identifies $Q_m$
as a primary Hopf surface. The reduced inverse image $P_m$ of
$E_0$ is smooth by (C.ER.7), since $E_0$ meets $R$ transversely.
Its map $p:P_m\to E_0$ has smooth elliptic fibres off $[0]$ and
one full fibre $mJ$ over $[0]$. The reduced fibre is
\[
                    J=\C^*/(q^m)^{\Z}.
                                                        \tag{C.ER.10}
\]
For example the isomorphism $E/\boldsymbol\mu_m\to J$ in
(C.ER.7) is $[w]\mapsto[w^m]$, which also specifies the finite
covering between the two elliptic curves.

Near $E_0\cap R$, choose coordinates with $R=(x=0)$,
$E_0=(y=0)$ and base map $t=xy$. In the smooth chart before
the free quotient the equation is $t=z^m y$; in (C.ER.4) it is
\[
                         t=a^m b y,\qquad b\ne0.
                                                        \tag{C.ER.11}
\]
It follows that the full central fibre is exactly $P_m+mQ_m$,
and its reduction has the two transverse components stated in
(C.ER.1). This calculation also retains the ramification at $J$.

We compute the canonical line without a classification assertion.
Choose a nowhere-zero holomorphic one-form $\omega$ on $E_0$.
On a local cyclic chart for $P_m$ as in (C.ER.7), the form
\[
                         p^*\omega\wedge\frac{dw}{w}
                                                        \tag{C.ER.12}
\]
is invariant under both $w\mapsto qw$ and the finite action.
Changing a local frame changes $dw/w$ by a form from the
one-dimensional base, which has zero wedge with $p^*\omega$.
It therefore glues to a global holomorphic two-form on $P_m$.
If $\omega=h(x)\,dx$ near $[0]$, its pullback is
\[
 m h(z^m)z^{m-1}\,dz\wedge\frac{dw}{w}.
                                                        \tag{C.ER.13}
\]
It has precisely the divisor $(m-1)J$, proving the canonical
isomorphism in (C.ER.2).

For every integer $k\ge0$ there is a precise pushforward equality
\[
        p_*\OO_{P_m}(kJ)
          =\OO_E\bigl(\lfloor k/m\rfloor[0]\bigr).
                                                        \tag{C.ER.14}
\]
A section on the left is a meromorphic function whose only allowed
poles are along $J$. Off $[0]$ it is holomorphic on every compact
elliptic fibre, so it is constant on that fibre and descends to a
holomorphic function on the base. On the cyclic chart at $[0]$,
its Laurent coefficients in $z$ are holomorphic functions on the
compact elliptic curve and hence constant. Finite-group invariance
forces the powers of $z$ to be multiples of $m$. The pole bound
$k$ is consequently exactly a pole bound $\lfloor k/m\rfloor$
in $x=z^m$. This proves both inclusions in (C.ER.14).
Applying it to $k=r(m-1)$ and applying Riemann--Roch on the elliptic
curve proves the displayed plurigenera. Their growth is linear
with positive slope $(m-1)/m$, so $\kappa(P_m)=1$.

Here is also an integral topological check that this surface is
non-K\"ahler. Choose a Hermitian metric on $L|_{E_0}$ and put
\[
 M_m=\{(v,\ell,a)\in W_m|_{E_0}:\|\ell\|=1\}.
                                                        \tag{C.ER.15}
\]
The maps $\ell=e^{mu}\ell_0$, $a=e^{-u}a_0$ identify
$W_m|_{E_0}=M_m\times\R$ smoothly. The generator (C.ER.5)
translates $u$ by $\log q$, so
$P_m\simeq M_m\times S^1$ smoothly. Over the complement of a
small disc about $[0]$, the canonical section trivializes $M_m$
as the product with a circle, whose coordinate is the phase of $a$.
On the disc choose the metric to be constant in the frame used
above. Then $M_m$ has the solid-torus coordinates
$a=\rho e^{\mathrm i\theta}$, $b=e^{\mathrm i\varphi}$, and
\[
              x=a^m b,\qquad \arg x=m\theta+\varphi.
                                                        \tag{C.ER.16}
\]
Let $\delta$ denote the boundary of the punctured elliptic base
and let $h$ denote its product circle. The meridian of this solid
torus maps to $m\delta+h$; its longitude maps to $\delta$.
The class $\delta$ is zero in the first homology of the once-punctured
elliptic curve. Mayer--Vietoris therefore kills $h$ and the solid-torus
longitude, leaving exactly the two free base classes. All spaces
and their boundary torus are connected, so the degree-zero part
adds no further class. Hence
$H_1(M_m,\Z)=\Z^2$ and $H_1(P_m,\Z)=\Z^3$.
The odd first Betti number excludes a K\"ahler metric, by Hodge
decomposition on a compact K\"ahler manifold.

The Euler calculation retains the multiple fibre. Over
$E_0\setminus\{[0]\}$, $P_m$ is locally a product with an elliptic
curve; its remaining reduced fibre $J$ is also elliptic.
Stratified Euler additivity gives $e(P_m)=0$.
The Hopf surface in (C.ER.9) is a principal elliptic bundle over
$\PP^1$, so $e(Q_m)=0$. Finally $e(J)=0$, and consequently
$e(P_m\cup_JQ_m)=0$.

Every map $\PP^1\to P_m$ is constant: it first maps constantly
to the elliptic base $E_0$, and then constantly to the reduced
elliptic fibre containing its image. Every map $\PP^1\to Q_m$
lifts to $\C^2\setminus\{0\}$ and is constant there by its
coordinate functions. Thus neither component is covered by
rational curves. The inverse-map resolution in (C.UC.1) applies
and proves that neither can be contracted by a bimeromorphic
comparison of smooth fillings. The Kodaira dimension of $Q_m$
is $-\infty$: a lifted pluricanonical form extends across the
origin by Hartogs, and its invariance under the contraction
$q^m\operatorname{id}$ forces all its power-series coefficients
to vanish. It cannot correspond birationally to $P_{m'}$.
Therefore any bimeromorphic comparison $U_m\dashrightarrow U_{m'}$
would identify $P_m$ birationally with $P_{m'}$.
Plurigenera are unchanged by a point blowup of a smooth compact
surface: a pluricanonical form descends off the centre and its
coefficients extend across that point by Hartogs. Resolution of
a surface bimeromorphism proves their birational invariance.
The different slopes $(m-1)/m$ in (C.ER.2) force $m=m'$.

For the persistence assertion, take a finite base change and choose
a local parameter in which it is $t=s^d$; a unit is absorbed by
its local $d$th root. At every point of the nonempty open subset
$P_m\setminus J$, the map $f_m$ is a submersion. Its base change
is the smooth chart $t=s^d$ with the two coordinates along $P_m$
unchanged. Thus a unique central component maps with degree one
onto this open subset. Its closure, followed by normalization and
any resolution, is birational to $P_m$. The preceding
non-uniruled-divisor comparison prevents its contraction under
any further bimeromorphic modification between smooth fillings.
This proves the asserted survival without imposing a description
on the other components of the base change.

Finally replace $\Delta$ throughout (C.ER.3) by $\PP^1$ and use
the same exceptional divisor, line bundle, and canonical section.
The local proofs of smoothness and properness still apply. Since
the base $V$ is now compact, the quotient is a compact smooth
threefold, with the same central fibre and with every other fibre
the product $E\times E$. This gives the claimed compact realizations.
The exact quotient-chart, canonical, homology, and plurigenus
calculations are retained in \path{code/E_elliptic_roots.py} and
its output.
\end{proof}

\paragraph{Literature comparison and an explicit source audit.}
The occurrence of a Kodaira-dimension-one component in a degeneration
of abelian surfaces already appears in
\cite[Example~3.4(1), p.~99]{Fujimoto1988Log}.
The preceding statement supplies its own quotient construction,
plurigenus distinctions, and finite-base-change comparison; it makes
no priority claim for the occurrence itself.

There is a precise reason not to use the divisibility assertion in
\cite[Theorem~2.8 and Corollary~2.10, pp.~96--97]{Fujimoto1988Log}
as an obstruction to (C.ER.4). For every $m\ge2$ consider the
diagonal Hopf threefold and the displayed map
\[
 \begin{split}
 H_m&=(\C^3\setminus\{0\})/
       \langle(a,b,c)\longmapsto(qa,q^m b,q^m c)\rangle,\\
 \Phi_m:H_m&\longrightarrow\PP^2,
       \qquad[a,b,c]\longmapsto[a^m:b:c].
 \end{split}                                             \tag{C.ER.17}
\]
The action is free and properly discontinuous. A compact annulus
between a small sphere and its image under the contraction meets every
orbit, so $H_m$ is compact and smooth. The three projective coordinates
in (C.ER.17) are multiplied by the same factor $q^m$, proving descent.
The map is holomorphic, surjective, and proper. Its fibre over
$[1:Y:Z]$ is parametrized by
$a\mapsto(a,Ya^m,Za^m)$ and is $\C^*/q^{\Z}$; its fibre over
$[0:Y:Z]$ is parametrized by
$s\mapsto(0,sY,sZ)$ and is $\C^*/(q^m)^{\Z}$.
In particular every fibre is connected and elliptic. These are all
the hypotheses in the definition of an elliptic threefold in
\cite[pp.~91--92]{Fujimoto1988Log}.
Near a point of the line $B=(X=0)$, use a chart with $b\ne0$ or
$c\ne0$. In the first chart $X/Y=a^m/b$, where $b$ is a unit.
It follows that
\[
             \Phi_m^*B=m\{a=0\},\qquad B^2=1,
             \qquad\pi_1(\PP^2\setminus B)=0.             \tag{C.ER.18}
\]
Thus (C.ER.17) is an explicit counterexample to both the asserted
divisibility and the asserted surjection to $\Z/m$ in that theorem
as printed. This conclusion follows from the verified map, rather
than from an inference that an incomplete proof makes a statement
false.

The problematic covering step can also be located exactly. For an
effective divisor $D$ with canonical section $s_D$, the construction
in that proof takes the equation $\xi^m=s_D^m$ in the total space of
$\OO(D)$. Its normalization is the disjoint union of the $m$ global
graphs $\xi=\zeta s_D$, $\zeta^m=1$, each isomorphic to the original
space. It therefore supplies no connected branched covering of the
base. Finally, (C.ER.17) is itself a special case of (C.ER.4): take
$V=\PP^2$, $L=\OO(1)$ and $\sigma=X$. A nonzero functional
$\ell$ on a projective line determines its unique vector $v$ with
$\ell(v)=1$; the equation $\sigma=a^m\ell$ becomes
$v=(a^m,b,c)$. Under this identification (C.ER.5) is the inverse
of the contraction in (C.ER.17), giving an actual isomorphism of
the two quotient constructions.


% C.2.UN.33 from part_C2_elliptic_types.tex

